\section{Image Processing and Comparison Method}\label{sec:image-processing}

\subsection{Computer-Vision Pipeline}

The image-processing pipeline will convert recorded Taylor-cone images into quantitative interface coordinates. The planned workflow is:
\begin{enumerate}
    \item select stable frames from video sequences,
    \item crop the region of interest around the meniscus,
    \item calibrate pixels to physical length using a known needle/nozzle dimension or calibration target,
    \item apply thresholding and/or edge detection,
    \item extract the cone boundary contour,
    \item fit local interface slopes near the cone region,
    \item compute the experimental half-angle,
    \item export $(r_{\mathrm{exp}},z_{\mathrm{exp}})$ profile coordinates and uncertainty estimates.
\end{enumerate}

Before physical data are collected, the pipeline will be tested on synthetic cone images with known ground-truth angles. The current end-to-end implementation has not passed this qualification: a diagnostic rerun on a nominal $49.29^{\circ}$ cone produced large angle and profile errors even at zero added noise. Consequently, this section is a proposed qualification protocol only; it is not evidence that the pipeline, camera calibration, or physical solver is validated. The repair gate requires held-out rendered images, explicit failure counts, fixed coordinate conventions, and predeclared thresholds before any experimental accuracy claim is made.

\subsubsection{Frame Selection}

Video sequences recorded during voltage ramps will contain hundreds to thousands of frames. Frame selection identifies the subset of frames that satisfy the stability criteria defined in Section~7. The selection algorithm will:
\begin{enumerate}
    \item Divide the video into segments corresponding to each voltage step (using timestamps or manual annotation).
    \item For each segment, compute a time series of shape metrics (e.g., apex position, interface curvature, centroid location).
    \item Identify the longest contiguous sub-interval within the segment where shape metrics remain within the stability tolerance bands.
    \item Select representative frames from the stable interval. Extract each frame separately and report the within-interval variation; do not treat adjacent video frames as independent experimental trials.
    \item Use a predeclared summary (for example, the median profile) only after checking that frame-to-frame variation is small relative to the measurement uncertainty.
\end{enumerate}

If no stable interval exists within a voltage segment, that segment will be classified as ``unstable'' and excluded from quantitative comparison.

\subsubsection{Calibration}

Pixel-to-millimeter calibration will use one of two methods:
\begin{enumerate}
    \item \textbf{Calibration target method:} A precision-machined scale or printed checkerboard pattern will be placed in the focal plane before the experiment. The pixel spacing corresponding to a known physical distance will be measured using edge detection or corner detection. This method is preferred when the camera-to-object distance is fixed and repeatable.
    \item \textbf{Needle diameter method:} The outer diameter of the needle will be measured independently using digital calipers or microscopy (e.g., $d_{\mathrm{out}} = 0.80 \pm 0.02$ for a standard 21-gauge needle). In the image, the needle diameter in pixels will be measured by fitting two parallel lines to the needle sidewalls far from the meniscus. The calibration scale is then $\text{mm/pixel} = d_{\mathrm{out}} / d_{\text{pixels}}$.
\end{enumerate}

Calibration uncertainty will be propagated through to all extracted dimensions. If lens distortion is significant, a distortion correction map will be applied using the calibrated-camera procedure described by Zhang~\cite{zhang2000calibration}. The distortion threshold, if one is used, will be declared from the calibration residuals rather than treated as a universal 1\% rule.

\subsubsection{Image Preprocessing and Edge Detection}

Raw grayscale images will be preprocessed to enhance the liquid-gas interface contrast:
\begin{enumerate}
    \item \textbf{Cropping:} Extract a region of interest (ROI) containing the needle tip and meniscus, discarding background pixels far from the interface.
    \item \textbf{Histogram equalization or normalization:} Adjust intensity distribution to maximize dynamic range within the ROI.
    \item \textbf{Gaussian blur (optional):} Apply a small Gaussian kernel (e.g., $\sigma = 1$--2 pixels) to reduce sensor noise, with kernel size chosen to preserve edge sharpness.
\end{enumerate}

Edge detection will use one of the following methods, selected based on performance on synthetic validation images:
\begin{itemize}
    \item \textbf{Canny edge detection:} A two-threshold gradient-based method that produces thin, connected edges. Hysteresis thresholds will be tuned on synthetic images and held fixed for all experimental trials.
    \item \textbf{Otsu thresholding + contour extraction:} Automatically select an intensity threshold that separates liquid (dark) from gas (bright) regions, then extract the boundary contour using a connected-component algorithm. This method is robust to lighting variation but may struggle with low-contrast or noisy images.
    \item \textbf{Sobel gradient magnitude + non-maximum suppression:} Compute the image gradient, threshold the magnitude, and suppress non-local-maximum pixels to produce thin edges. This is a simpler alternative to Canny.
\end{itemize}

The chosen method will be documented in the final validation report along with all parameter values (e.g., Canny thresholds, Gaussian $\sigma$, Otsu method).

\subsubsection{Contour Extraction and Smoothing}

After edge detection, the interface contour will be extracted as an ordered list of pixel coordinates $\{(x_i, y_i)\}_{i=1}^{N_{\text{contour}}}$. Post-processing steps include:
\begin{enumerate}
    \item \textbf{Contour selection:} If multiple contours are detected (e.g., needle edge, interface, artifacts), select the contour that:
    \begin{itemize}
        \item originates near the needle tip,
        \item extends into the region between the needle and extractor,
        \item has monotonically increasing axial coordinate (no backtracking),
        \item has the longest arclength among candidates.
    \end{itemize}
    \item \textbf{Axisymmetry enforcement:} Since the experiment observes a 2D projection of an axisymmetric interface, the pipeline will extract one side of the contour (e.g., left or right boundary) and assume radial symmetry. If both sides are visible, the two half-profiles can be averaged to reduce asymmetry artifacts from optical misalignment.
    \item \textbf{Smoothing (optional):} Apply a low-pass filter (e.g., Savitzky--Golay filter, cubic spline with smoothing parameter) to reduce pixel-level noise while preserving the overall cone shape. The smoothing strength will be validated on synthetic images to ensure it does not bias the extracted angle.
\end{enumerate}

\subsubsection{Coordinate Transformation to $(r, z)$}

The pixel coordinates will be transformed into the solver's axisymmetric coordinate system:
\begin{enumerate}
    \item Apply the calibration scale: $(x_{\text{mm}}, y_{\text{mm}}) = \text{scale} \times (x_{\text{px}}, y_{\text{px}})$.
    \item Identify the needle tip or apex as the coordinate origin. For needle-tip referencing, locate the point where the interface detaches from the needle (sharp curvature change or minimum radius). For apex referencing, locate the point of maximum axial penetration into the gas.
    \item Define the axial direction $z$ (parallel to needle axis, positive toward extractor) and the radial direction $r$ (perpendicular to axis, positive away from centerline).
    \item Transform the contour into $(r, z)$ coordinates: $r_i = |x_i|$ (distance from axis), $z_i = y_i$ (axial position).
\end{enumerate}

The experimental profile is now represented as $\{(r_i^{\exp}, z_i^{\exp})\}_{i=1}^{N_{\exp}}$.

\subsubsection{Half-Angle Measurement}

The cone half-angle will be extracted by fitting a straight line to the flank region of the experimental profile. The fitting procedure is:
\begin{enumerate}
    \item \textbf{Define flank region:} Select the interval $z \in [z_{\min}, z_{\max}]$ where the interface is approximately conical. Exclude the apex region (high curvature, rounded cap) and the base region (near needle tip, attachment artifacts). A typical choice is the middle 50\% of the visible cone length, e.g., $z_{\min} = z_{\text{apex}} - 0.25 L_{\text{cone}}$, $z_{\max} = z_{\text{apex}} - 0.75 L_{\text{cone}}$, where $L_{\text{cone}}$ is the total cone length from apex to needle.
    \item \textbf{Least-squares line fit:} Fit $r = m z + b$ to the flank points using weighted or unweighted least squares. The slope $m = \tan(\alpha)$ gives the half-angle:
    \begin{equation}
    \alpha_{\exp} = \arctan(m).
    \end{equation}
    \item \textbf{Uncertainty estimate:} The standard error of the slope $\sigma_m$ from the regression can be propagated to the angle:
    \begin{equation}
    \sigma_{\alpha} = \frac{\sigma_m}{1 + m^2}.
    \end{equation}
    If multiple frames are analyzed, the angle standard deviation across frames provides an additional uncertainty estimate.
\end{enumerate}

The flank-window selection will be validated on synthetic images to ensure it produces accurate angles for known cone geometries and is not biased by the window choice.

\subsection{Profile Alignment}

Simulation and experiment will be compared in the same coordinate system. Experimental pixel coordinates will be converted to millimeters and aligned so that the extracted apex or another reproducible reference point matches the simulation coordinate origin. The experimental radius $r_{\mathrm{exp}}(z_i)$ will be interpolated onto the simulation $z_i$ locations, or both profiles will be interpolated onto a common coordinate grid.

\paragraph{Alignment Strategy.} Three alignment methods will be considered:

\begin{enumerate}
    \item \textbf{Apex alignment:} Translate the experimental profile so that its apex (point of maximum $z$ or minimum $r$) coincides with the simulated apex. This is geometrically natural but may be sensitive to apex-region noise or cap artifacts.
    \item \textbf{Needle-tip alignment:} Translate both profiles so that the needle-tip detachment point (experimentally measured, simulated as the cone base) is at the origin. This is robust to apex uncertainty but requires clear identification of the detachment point in both experiment and simulation.
    \item \textbf{Flank-region least-squares alignment:} Minimize the profile RMSE over translations and small rotations, using only the flank region (excluding apex and base). This is robust to localized artifacts but adds degrees of freedom that could artificially reduce error.
\end{enumerate}

The primary method will be apex alignment unless preliminary validation on synthetic images shows it is unreliable. The alignment choice will be frozen before processing experimental data.

\paragraph{Interpolation.} After alignment, one profile will be interpolated onto the other's coordinate grid. Linear or cubic spline interpolation will be used. If the experimental profile has higher spatial sampling than the simulation grid, the experimental points will be interpolated onto the simulation $z_i$ locations. If the simulation is finer, the opposite will be done. Points where one profile extends beyond the other (e.g., experimental profile reaches closer to the base) will be excluded from RMSE calculation to avoid extrapolation artifacts.

\subsection{Validation Metrics}

The primary profile metric will be the root mean square error:
\begin{equation}
\mathrm{RMSE}_{R}=\sqrt{\frac{1}{N}\sum_{i=1}^{N}\left(R_{\mathrm{sim}}(z_i)-R_{\mathrm{exp}}(z_i)\right)^2}.
\end{equation}
A normalized version will divide by a characteristic length $L_c$, such as electrode spacing or observed cone length:
\begin{equation}
\mathrm{nRMSE}_{R}=\frac{\mathrm{RMSE}_{R}}{L_c}.
\end{equation}
The cone-angle error will be
\begin{equation}
\Delta\alpha=|\alpha_{\mathrm{sim}}-\alpha_{\mathrm{exp}}|.
\end{equation}
For onset voltage,
\begin{equation}
\epsilon_V=\frac{|V_{\mathrm{sim}}-V_{\mathrm{exp}}|}{V_{\mathrm{exp}}}.
\end{equation}
Repeated trials will report mean, standard deviation, and uncertainty intervals where possible.

\paragraph{Profile RMSE Decomposition.} To diagnose the source of profile mismatch, the radial error can be decomposed into components:
\begin{itemize}
    \item \textbf{Angle-driven error:} If the simulated and experimental half-angles differ by $\Delta\alpha$, the radial profiles will diverge as $\Delta r \approx z \tan(\Delta\alpha)$ along the flank. This component can be estimated by comparing the RMSE from the actual profiles to the RMSE from two cones with half-angles $\alpha_{\mathrm{sim}}$ and $\alpha_{\mathrm{exp}}$ but otherwise identical geometry.
    \item \textbf{Apex-region error:} Near the apex, the experimental profile may have a finite-radius cap due to finite resolution or physical regularization (surface tension smoothing, finite needle radius), while the simulation may impose a sharper or differently sized cap. The apex-region contribution can be isolated by computing RMSE separately for $z < z_{\text{apex}} - \delta$ (apex window) and $z > z_{\text{apex}} - \delta$ (flank window).
    \item \textbf{Base-region error:} Near the needle attachment, contact-line dynamics, wetting behavior, and needle geometry artifacts may cause deviations not captured by the idealized cone model. Base-region RMSE can be computed separately.
\end{itemize}

\paragraph{Statistical Significance.} Repeated trials at the same condition will produce a distribution of measured angles, onset voltages, and profile RMSEs. Agreement between simulation and experiment will be assessed by comparing the simulation prediction to the experimental mean $\pm$ standard deviation (or standard error of the mean if the trial count is small). If the prediction falls within the experimental uncertainty band, agreement is strong. If not, the discrepancy magnitude will be reported as a multiple of the experimental standard deviation (e.g., ``simulation predicts $\alpha = 47.5^{\circ}$, experiment measures $49.2 \pm 0.8^{\circ}$, discrepancy is $2.1 \sigma$'').

\subsection{Synthetic Validation of the Image Processing Pipeline}

Before analyzing experimental data, the entire image-processing pipeline will be validated on synthetic images generated from the solver output. The validation procedure is:

\paragraph{Step 1: Generate Synthetic Images.}
\begin{enumerate}
    \item Run the solver at a known geometry (e.g., electrode spacing $L = 10$, needle diameter $d = 0.8$, surface tension $\gamma = 0.022$, applied voltage $V = 3$).
    \item Extract the predicted interface profile $R_{\mathrm{sim}}(z)$ and half-angle $\alpha_{\mathrm{sim}}$ (ground truth).
    \item Render the profile as a 2D silhouette image using a graphics library (e.g., matplotlib filled contour, PIL polygon rendering, or OpenCV drawing functions).
    \item Embed the silhouette in a realistic image background: add a needle representation, set pixel resolution (e.g., 1024$\times$768), apply calibration scale (e.g., 50), render a backlight gradient.
    \item Corrupt the synthetic image with realistic noise and blur: add Gaussian intensity noise (SNR 40--60~dB), apply Gaussian blur ($\sigma = 1$--2~pixels), simulate lens distortion if applicable, add pixel quantization (8-bit or 12-bit depth).
\end{enumerate}

\paragraph{Step 2: Process Synthetic Images Through Pipeline.}
\begin{enumerate}
    \item Apply the full image-processing pipeline (edge detection, contour extraction, calibration, angle measurement) to the synthetic image.
    \item Extract $\alpha_{\mathrm{extracted}}$ and $R_{\mathrm{extracted}}(z)$ from the pipeline output.
    \item Compute the round-trip errors:
    \begin{equation}
    \Delta\alpha_{\mathrm{synthetic}} = |\alpha_{\mathrm{extracted}} - \alpha_{\mathrm{sim}}|,
    \end{equation}
    \begin{equation}
    \mathrm{RMSE}_{\mathrm{synthetic}} = \sqrt{\frac{1}{N}\sum_{i=1}^{N}\left(R_{\mathrm{extracted}}(z_i) - R_{\mathrm{sim}}(z_i)\right)^2}.
    \end{equation}
\end{enumerate}

\paragraph{Step 3: Acceptance Criteria.}
The image-processing pipeline will be accepted as valid if the synthetic validation shows:
\begin{itemize}
    \item $\Delta\alpha_{\mathrm{synthetic}} < 1^{\circ}$ for cone angles in the range $40^{\circ}$ to $55^{\circ}$ under moderate noise (SNR $\geq 40$~dB),
    \item $\mathrm{nRMSE}_{\mathrm{synthetic}} < 2\%$ of the cone length,
    \item Round-trip error does not increase systematically with angle, cone length, or noise level (indicating method robustness).
\end{itemize}

If this future synthetic qualification fails its predeclared criteria, the edge-detection method, smoothing parameters, or flank-window selection will be revised before experimental images are analyzed. Until that gate passes, no synthetic image result is presented as evidence that the image-analysis method is unbiased, experimentally calibrated, or physically validated.

\subsection{Planned Output Figures}

The final paper will include: raw image frames, detected-edge overlays, extracted cone profiles, simulated-versus-experimental profile overlays, half-angle comparison plots, onset-voltage comparison plots, and runtime scaling figures.

\paragraph{Representative Image Sequence.} A figure will show 3--5 frames from a single voltage ramp: (a) no cone (hemispherical meniscus at low voltage), (b) deformed meniscus (subcritical), (c) stable Taylor cone at onset, (d) cone with incipient jet or spray (post-onset). Each frame will include the detected edge overlay and annotated voltage.

\paragraph{Profile Comparison.} For each validated trial, a plot will overlay the experimental profile $R_{\mathrm{exp}}(z)$ (solid line with error band) and the frozen simulation prediction $R_{\mathrm{sim}}(z)$ (dashed line) in the same coordinate system. The half-angles will be annotated with fitted flank lines. Residual error $R_{\mathrm{sim}}(z) - R_{\mathrm{exp}}(z)$ will be plotted as a function of $z$ in a separate panel.

\paragraph{Angle and Voltage Comparison.} Scatter plots or bar charts will compare simulated versus experimental half-angles and onset voltages across all trials. Error bars will represent experimental uncertainty (standard deviation across repeated trials). A diagonal line $y = x$ will indicate perfect agreement. The plots will quantify systematic bias (mean offset) and random scatter (RMSE about the diagonal).

\paragraph{Voltage-Sweep Landscape.} For at least one trial, a figure will show the experimental regime map (voltage vs qualitative state: no cone, unstable, stable, spray) overlaid with the simulated residual landscape (residual RMS vs voltage). This will illustrate how the numerical onset criterion aligns with the experimental onset observation.

\paragraph{Synthetic Validation Results.} A supplementary figure will document the synthetic validation: (a) example synthetic image with ground-truth profile overlaid, (b) extracted profile from pipeline, (c) angle error and profile RMSE as functions of noise level and cone angle, demonstrating pipeline accuracy and robustness.
